\documentclass[a4paper,12pt]{article}
%====================== PACKAGES ======================
\usepackage[english]{babel}
\usepackage[utf8]{inputenc}
\usepackage{float}
\usepackage{amsmath}
\usepackage{graphicx}
\usepackage[colorinlistoftodos]{todonotes}
\usepackage{url}
\usepackage{hyperref}
\usepackage{array}
\usepackage{tabularx}
\usepackage{setspace}
\usepackage{abstract}
\usepackage[T1]{fontenc}
\usepackage[top=2cm, bottom=2cm, left=2cm, right=2cm]{geometry}
\usepackage{subfig}
\usepackage{newunicodechar}
\newunicodechar{₂}{$_2$}
\usepackage{amsmath}
\usepackage{amsfonts}
\usepackage{amssymb}
\usepackage{hyperref}
\usepackage[backend=biber,style=numeric]{biblatex}
\usepackage{lipsum}
\usepackage{xcolor}
\usepackage{caption}
\usepackage{amsfonts}
\usepackage{amsmath, amssymb}
\usepackage{tikz, pgfplots}
\pgfplotsset{compat=1.18}
\usepackage{booktabs}
\usepackage{multirow}
\usepackage{siunitx}
\sisetup{table-format=2.3, table-number-alignment=center}

%====================== INFORMATION ET REGLES ======================
\setcounter{secnumdepth}{4}
\setcounter{tocdepth}{4}
\newcommand{\HRule}{\rule{\linewidth}{0.5mm}}

\begin{document}

%====================== TITLE PAGE ======================
\begin{titlepage}
\begin{center}
\includegraphics[scale=0.15]{./logo-epfl}~\\[1cm]
\textsc{\Large }\\[0.5cm]
\HRule \\[0.4cm]
{\huge \bfseries Investments Project \\[0.4cm] }
{\large \bfseries Investments \\[0.4cm] }
\HRule \\[1.5cm]
\vspace{3cm}
\begin{minipage}{0.4\textwidth}
\begin{center} \large
\emph{Group 14:}\\
    Alexandre \textsc{Vallet} \\
    Leonard \textsc{Rambeau} \\
    Victor \textsc{Legrand} \\
    Loric \textsc{Rey}
\end{center}
\end{minipage}
\vfill
{\large \today}
\end{center}
\end{titlepage}

%====================== TABLE OF CONTENTS ======================
\tableofcontents
\newpage

%====================== SECTION 1: INTRODUCTION ======================
\section{Introduction}

This report studies currency excess returns and systematic trading strategies across the G10
foreign exchange (FX) market.

The core objective is twofold. First, to characterise the statistical properties of individual
currency excess returns. Second, to evaluate three well-established currency factor strategies
--- carry, momentum, and dollar --- and to assess whether their return premia are statistically
robust, economically significant, and mutually independent.

We then move from individual factors to portfolio construction: in-sample mean-variance
portfolios in Section~\ref{sec:insample}, and out-of-sample conditional mean-variance
efficient (CMVE) portfolios in Section~\ref{sec:oos}.

\paragraph{Why FX?}
Currency markets offer a distinct testing ground for asset pricing. Unlike equities, FX
excess returns are zero-sum across investors and tightly linked to interest rate differentials
and exchange rate dynamics. The persistence of carry premia, in particular, constitutes a
well-documented violation of Uncovered Interest Rate Parity (UIP) and lies at the heart of
the ``forward premium puzzle''.

\paragraph{Sample.}
We study six G10 currencies against the US dollar (AUD, CAD, EUR, GBP, JPY, NZD) over
the period \textbf{January 1995 to December 2024}, with December 1990 as the earliest data
point for signal construction and a holdout window from January 1991 through December 1994
used to initialise out-of-sample estimators. All results are reported on a monthly basis
and annualised where appropriate.

%====================== SECTION 2: DATA ======================
\section{Data}
\label{sec:data}

\subsection{FX Rates}

Spot and one-month forward mid-quote exchange rates are sourced from WRDS Datastream
(\texttt{tr\_ds\_equities.ds2fxrate}), which stores FX series in its equities database.

We use the \texttt{exrateintcode} identifiers specified in the project brief
(e.g.\ 2427 for AUD spot, 2601 for AUD forward) for the six currency pairs. Data are
available at daily frequency from December 1990 to December 2024. To obtain a monthly
panel, we take the last daily observation within each calendar month.

All rates are expressed in USD per unit of foreign currency (USD/FCU); series originally
quoted in FCU/USD (e.g.\ CAD, JPY) are inverted, and series with inconsistent quoting
conventions across spot and forward legs are handled on a per-series basis.

\subsection{Interest Rates}

Monthly short-term interest rates are downloaded from FRED using the OECD MEI series
\texttt{IRSTCI01}, which reports call money / overnight interbank rates, annualised,
not seasonally adjusted, in percentage points per year.

The seven series used are:
\begin{itemize}
    \item USD (\texttt{IRSTCI01USM156N}), AUD (\texttt{IRSTCI01AUM156N}),
          CAD (\texttt{IRSTCI01CAM156N}),
    \item GBP (\texttt{IRSTCI01GBM156N}), JPY (\texttt{IRSTCI01JPM156N}),
          NZD (\texttt{IRSTCI01NZM156N}),
    \item a spliced EUR series using the German DEM rate
          (\texttt{IRSTCI01DEM156N}) before January 1999 and the Euro Area rate
          (\texttt{IRSTCI01EZM156N}) from January 1999 onward.
\end{itemize}

All rates are converted from annualised percentage points to monthly decimals by dividing
by 1200, as required by the excess return formula.

\subsection{Final Panel}

The merged panel contains 409 monthly observations from December 1990 to December 2024,
with 19 columns: spot rates $S^i_t$, one-month forward rates $F^i_t$, and risk-free rates
$r^i_t$ for the six foreign currencies, plus the USD rate $r^{US}_t$. The panel is
complete with no missing values.

%====================== SECTION 3: EXCESS RETURNS ======================
\section{Currency Excess Returns}
\label{sec:excess}

\subsection{Construction}

The excess return in USD from a \$1 investment in currency $i$ during month $t+1$ is:

\begin{equation}
    X^i_{t+1} = \frac{S^i_{t+1}}{S^i_t}\left(1 + r^i_t\right) - \left(1 + r^{US}_t\right)
    \label{eq:excess}
\end{equation}

This represents the payoff from borrowing \$1 at the US rate $r^{US}_t$, converting into
$1/S^i_t$ units of foreign currency, investing at the foreign rate $r^i_t$, and
reconverting at the realised spot rate $S^i_{t+1}$ one month later.

The excess return thus combines two components:
the \textbf{interest rate differential} $(r^i_t - r^{US}_t)$
and the \textbf{exchange rate change} $(S^i_{t+1}/S^i_t - 1)$.

Excess returns are computed over the full sample January 1991 to December 2024
(408 observations). All performance statistics below are reported on the
\textbf{evaluation sample}: January 1995 to December 2024 (360 observations).

\subsection{Summary Statistics --- Table 1}

Table~\ref{tab:table1} reports summary statistics for the six individual currency excess
returns over the evaluation sample. Currencies are ordered alphabetically.

\begin{table}[H]
\centering
\caption{Currency Excess Returns --- Summary Statistics (Jan 1995 -- Dec 2024)}
\label{tab:table1}
\renewcommand{\arraystretch}{1.3}
\begin{tabular}{lrrrrrrrr}
\toprule
 & \textbf{AUD} & \textbf{CAD} & \textbf{EUR} & \textbf{GBP} & \textbf{JPY} & \textbf{NZD} \\
\midrule
Mean (\%, ann.)        & 1.408    & 0.334    & $-$0.888 & 0.101    & $-$3.302 & 2.258  \\
$t$-statistic          & 0.656    & 0.227    & $-$0.527 & 0.067    & $-$1.685 & 1.013  \\
Std (\%, ann.)         & 11.758   & 8.064    & 9.222    & 8.278    & 10.733   & 12.210 \\
Sharpe ratio           & 0.120    & 0.041    & $-$0.096 & 0.012    & $-$0.308 & 0.185  \\
Skewness               & $-$0.192 & $-$0.281 & 0.035    & $-$0.206 & 0.658    & $-$0.111 \\
Excess kurtosis        & 1.303    & 2.912    & 1.214    & 1.222    & 3.254    & 1.242  \\
$N$ obs                & 360      & 360      & 360      & 360      & 360      & 360    \\
$\bar{r}^i$ (\%, ann.) & 3.968    & 2.586    & 1.750    & 3.001    & 0.140    & 4.459  \\
\bottomrule
\end{tabular}
\caption*{\footnotesize Mean and standard deviation are annualised (multiply monthly by 12 and $\sqrt{12}$
respectively). Sharpe ratio is annualised mean divided by annualised standard deviation.
$t$-statistic tests $H_0: \mathbb{E}[X^i] = 0$. Skewness and excess kurtosis are
bias-corrected. $\bar{r}^i$ is the average foreign risk-free rate over the evaluation sample.}
\end{table}

\subsection{Discussion}

Several patterns emerge from Table~\ref{tab:table1}.

\paragraph{Cross-sectional heterogeneity in mean returns.}
Mean excess returns vary substantially across currencies, ranging from $-3.30\%$ for JPY
to $+2.26\%$ for NZD.

This dispersion is closely aligned with the cross-section of foreign risk-free rates: AUD
($\bar{r}^i = 3.97\%$) and NZD ($\bar{r}^i = 4.46\%$) are the highest-yielding currencies
and earn the largest positive excess returns, while JPY ($\bar{r}^i = 0.14\%$) is the
lowest-yielding and earns the most negative excess return.

This pattern is qualitatively consistent with a carry premium: high-rate currencies tend
to outperform low-rate currencies on a risk-adjusted basis, violating the UIP prediction
that exchange rate changes should exactly offset rate differentials.

\paragraph{Statistical insignificance at the individual currency level.}
Despite the cross-sectional dispersion in means, no individual currency excess return is
statistically distinguishable from zero at conventional thresholds.

The highest $t$-statistic in absolute value belongs to JPY ($-1.69$), which still falls
short of the 1.96 threshold for 5\% significance. This reflects the high idiosyncratic
volatility of single-currency positions (annual standard deviations ranging from 8.1\% to
12.2\%) relative to their modest average returns.

Efficient FX investing therefore requires diversification across currencies and conditioning
on systematic signals, rather than individual currency bets.

\paragraph{Fat tails and skewness.}
Excess kurtosis is positive for all currencies, ranging from 1.21 (EUR) to 3.25 (JPY),
indicating heavier tails than a normal distribution.

The skewness pattern is noteworthy: JPY displays \emph{positive} skewness ($+0.658$),
reflecting the occasional sharp yen appreciations during global risk-off episodes (carry
unwind events, e.g.\ 2008, 2020). High-yield currencies AUD, CAD, and GBP display mild
negative skewness, consistent with the well-known crash risk faced by carry investors who
are long these currencies.

%====================== SECTION 4: CARRY ======================
\section{Carry Strategies}
\label{sec:carry}

\subsection{Motivation and Signal}

The carry trade exploits the empirical failure of UIP. Rather than depreciating by the
interest rate differential as UIP predicts, high-rate currencies tend to appreciate or
remain stable, generating excess returns for the carry investor.

The signal for carry strategies is the interest rate differential:
\begin{equation}
    \Delta r^i_t = r^i_t - r^{US}_t.
\end{equation}

\subsection{Strategy Construction}

\paragraph{CS-CARRY (Cross-Sectional Carry).}
Each month, currencies are ranked by $\Delta r^i_t$ (rank 1 = lowest differential, rank
$N = 6$ = highest). Portfolio weights are assigned as
\begin{equation}
    w^i_t = Z \left( \mathrm{Rank}^i_t - \frac{N+1}{2} \right),
\end{equation}
where $Z$ is chosen such that $\sum_{i:\, w^i_t > 0} w^i_t = 1$, making the strategy
self-financing with exactly \$1 long and \$1 short. The strategy bets on the
\emph{relative} performance of high-rate versus low-rate currencies within the
cross-section.

\paragraph{TS-CARRY (Time-Series Carry).}
Each currency is traded independently based on the sign of its own differential:
\begin{equation}
    w^i_t = \frac{\mathrm{sign}(\Delta r^i_t)}{N}.
\end{equation}
This strategy takes a \$1/N position in each currency independently of the others, going
long if the foreign rate exceeds the US rate and short otherwise.

All signals are formed at month-end $t$ and used to trade the following month's return
$X^i_{t+1}$, with no look-ahead bias.

\subsection{Results --- Table 2 (Carry Rows)}

\begin{table}[H]
\centering
\caption{Factor Strategies --- Performance Statistics (Jan 1995 -- Dec 2024)}
\label{tab:table2}
\renewcommand{\arraystretch}{1.3}
\begin{tabular}{lrrrrrr}
\toprule
 & \textbf{CS-CARRY} & \textbf{TS-CARRY} & \textbf{CS-MOM} & \textbf{TS-MOM} & \textbf{DOLLAR} & \textbf{DOLLAR-CARRY} \\
\midrule
Mean (\%, ann.)  & 3.691 & 2.772 & 0.235  & 1.564 & $-$0.015 & 2.128 \\
$t$-statistic    & 2.166 & 2.568 & 0.153  & 1.293 & $-$0.011 & 1.543 \\
Std (\%, ann.)   & 9.334 & 5.912 & 8.405  & 6.624 & 7.579    & 7.554 \\
Sharpe ratio     & 0.395 & 0.469 & 0.028  & 0.236 & $-$0.002 & 0.282 \\
\bottomrule
\end{tabular}
\caption*{\footnotesize All statistics computed over the evaluation sample (360 monthly observations).
Mean and standard deviation are annualised. $t$-statistic tests $H_0: \mathbb{E}[R] = 0$.}
\end{table}

\paragraph{Significance.}
Both carry strategies deliver statistically significant positive mean returns. CS-CARRY
earns $3.69\%$ per annum ($t = 2.17$) with a Sharpe ratio of $0.40$, while TS-CARRY earns
$2.77\%$ ($t = 2.57$) with a Sharpe ratio of $0.47$. Both $t$-statistics exceed the 1.96
threshold.

\paragraph{TS dominates CS in Sharpe terms.}
TS-CARRY achieves a \emph{higher} Sharpe ratio despite its lower mean, because its
volatility ($5.91\%$) is substantially lower than CS-CARRY's ($9.33\%$). This reflects
the fact that TS-CARRY positions are sized equally across currencies regardless of the
magnitude of their differential, which limits concentration risk.

\subsection{CS-CARRY Portfolio Weights Over Time}

Figure~\ref{fig:weights} plots the CS-CARRY portfolio weight on each currency over time.
The weights reveal strong persistence:
\begin{itemize}
    \item AUD and NZD are persistently long throughout the sample, consistent with their
          structurally high interest rates relative to the US.
    \item JPY is persistently short, reflecting its near-zero rate environment since the
          mid-1990s.
    \item CAD, EUR, and GBP occupy intermediate ranks and switch sides as their monetary
          policy cycles diverge from or converge with the US.
\end{itemize}

An exception occurs around 2012--2017, when the Federal Reserve's zero-interest-rate policy
(ZIRP) compressed the differential between the US rate and JPY, temporarily reducing JPY's
short position.

\begin{figure}[H]
    \centering
    \includegraphics[width=\textwidth]{fig_cscarry_weights.pdf}
    \caption{CS-CARRY portfolio weights on each of the six currencies (January 1995 -- December 2024).
    Dark fill = long position; light fill = short position. AUD and NZD are persistently long
    (high rates); JPY is persistently short (near-zero rate).}
    \label{fig:weights}
\end{figure}

\subsection{Forward Discount and Interest Rate Differential}

The net forward discount is defined (in logs) as
\begin{equation}
    fd^i_t = \log\frac{S^i_t}{F^i_t}.
\end{equation}

By Covered Interest Rate Parity (CIP), which holds very precisely in G10 markets via
forward market arbitrage:
\begin{equation}
    \frac{F^i_t}{S^i_t} = \frac{1 + r^{US}_t}{1 + r^i_t}
    \quad \Longrightarrow \quad
    fd^i_t = \log(1 + r^i_t) - \log(1 + r^{US}_t) \approx \Delta r^i_t.
\end{equation}

The log-linear approximation holds for small rates. Table~\ref{tab:corr} reports the
time-series correlation between $fd^i_t$ and $\Delta r^i_t$ for each currency.

\begin{table}[H]
\centering
\caption{Time-Series Correlation between Forward Discount $fd^i_t$ and Rate Differential $\Delta r^i_t$}
\label{tab:corr}
\renewcommand{\arraystretch}{1.3}
\begin{tabular}{lrrrrrr|r}
\toprule
 & AUD & CAD & EUR & GBP & JPY & NZD & \textbf{Average} \\
\midrule
$\mathrm{Corr}(fd^i_t, \Delta r^i_t)$ & 0.793 & 0.756 & 0.944 & 0.809 & 0.973 & 0.803 & \textbf{0.846} \\
\bottomrule
\end{tabular}
\caption*{\footnotesize Computed over the full sample (December 1990 -- December 2024).
Cross-sectional average reported in bold.}
\end{table}

The cross-sectional average correlation of $0.846$ is high and positive, consistent with
CIP theory.

Deviations from unity arise from two sources: the log-level approximation error (negligible
for typical G10 rate levels), and minor data-source timing differences between Datastream
forward quotes and FRED interest rate observations.

The sign is unambiguously positive: the carry signal embedded in the forward market
($fd^i_t$) and the interest rate differential ($\Delta r^i_t$) are two equivalent
representations of the same underlying information.

%====================== SECTION 5: MOMENTUM ======================
\section{Momentum Strategies}
\label{sec:momentum}

\subsection{Motivation and Signal}

Momentum strategies exploit the empirical tendency of recent currency appreciations to
persist. The signal is the 12-month spot appreciation against the dollar:
\begin{equation}
    s^i_{t-12,t} = \frac{S^i_t}{S^i_{t-12}} - 1.
\end{equation}

This signal is constructed on the full sample starting December 1991 (the first month for
which both $S^i_t$ and $S^i_{t-12}$ are available given our data begins December 1990).

The economic rationale rests on underreaction: markets are slow to incorporate information
about trend-following currency dynamics, leading to persistent price momentum.

\subsection{Strategy Construction}

\paragraph{CS-MOM (Cross-Sectional Momentum).}
Currencies are ranked each month by $s^i_{t-12,t}$ and assigned weights following the same
rank-weighting formula as CS-CARRY (Equation~2).

\paragraph{TS-MOM (Time-Series Momentum).}
Each currency is traded independently based on the sign of its own momentum signal, using
the same equal-exposure formula as TS-CARRY (Equation~3).

\subsection{Results --- Table 2 (Momentum Rows)}

Table~\ref{tab:table2} shows that FX momentum is substantially weaker than carry over this
sample.

\begin{itemize}
    \item \textbf{CS-MOM} earns only $0.24\%$ per annum ($t = 0.15$, Sharpe $= 0.03$).
    \item \textbf{TS-MOM} earns $1.56\%$ per annum ($t = 1.29$, Sharpe $= 0.24$).
\end{itemize}

\textbf{Neither strategy is statistically significant at conventional thresholds.} This is
consistent with the broader literature, which finds that FX momentum is less robust than
equity momentum and more concentrated in specific sub-periods of trending exchange rate
regimes (Menkhoff et al., 2012).

As with carry, TS-MOM dominates CS-MOM in Sharpe terms because its sign-based construction
produces lower volatility ($6.62\%$ vs $8.40\%$).

\subsection{Regression of TS-CARRY on Momentum Factors}

To assess whether momentum captures part of the carry premium, we estimate
\begin{equation}
    R^{TS\text{-}CARRY}_{t+1} = \alpha + \beta_1 R^{CS\text{-}MOM}_{t+1}
                                + \beta_2 R^{TS\text{-}MOM}_{t+1} + u_{t+1},
\end{equation}
using Newey-West HAC standard errors with 12 lags. Results appear in column~(1) of
Table~\ref{tab:reg}.

\begin{table}[H]
\centering
\caption{Regression of TS-CARRY on Momentum and Dollar Factors (HAC standard errors, 12 lags)}
\label{tab:reg}
\renewcommand{\arraystretch}{1.3}
\begin{tabular}{lcccc}
\toprule
 & \multicolumn{2}{c}{\textbf{(1) MOM only}} & \multicolumn{2}{c}{\textbf{(2) MOM + DOLLAR}} \\
\cmidrule(lr){2-3} \cmidrule(lr){4-5}
 & Coef. & $t$-stat & Coef. & $t$-stat \\
\midrule
Intercept ($\alpha$)     & 0.0023 & 2.79$^{**}$ & 0.0022 & 3.02$^{***}$ \\
CS-MOM ($\beta_1$)       & 0.0196 & 0.27        & 0.0193 & 0.30 \\
TS-MOM ($\beta_2$)       & 0.0013 & 0.01        & 0.0802 & 0.63 \\
DOLLAR ($\beta_3$)       & ---    & ---         & 0.3195 & 2.49$^{**}$ \\
\midrule
$R^2$                    & \multicolumn{2}{c}{0.08\%} & \multicolumn{2}{c}{16.1\%} \\
Standalone $t(\bar{R}^{TS\text{-}CARRY})$ & \multicolumn{4}{c}{2.57} \\
\bottomrule
\end{tabular}
\caption*{\footnotesize $^{**}$ significant at 5\%; $^{***}$ significant at 1\%.
$N = 360$ monthly observations. HAC standard errors (Newey-West, 12 lags).
Intercepts are monthly; annualised alpha = $\alpha \times 12 \times 100$.}
\end{table}

\paragraph{Carry and momentum are nearly orthogonal.}
In the MOM-only model, both momentum coefficients are statistically indistinguishable from
zero ($t < 0.30$) and $R^2 = 0.08\%$, confirming that carry and momentum capture different
information in expected return space.

\paragraph{The intercept t-stat rises through the variance channel, not the mean.}
The alpha $t$-statistic increases from 2.57 (standalone TS-CARRY mean) to 2.79 (after MOM
controls). Since the beta coefficients on momentum are positive but near zero, the alpha
itself is essentially unchanged.

The mild increase in $t(\alpha)$ therefore reflects a small reduction in residual variance,
not an increase in the carry mean. In short: adding momentum to a carry portfolio
\emph{marginally} improves the precision of carry estimation, but momentum does not subsume
or amplify carry.

%====================== SECTION 6: DOLLAR ======================
\section{Dollar Strategies}
\label{sec:dollar}

\subsection{Motivation and Signal}

Carry and momentum strategies are dollar-neutral by construction: they long some currencies
and short others, so any broad USD movement cancels out.

Dollar strategies instead take a \emph{common directional position} on all foreign
currencies simultaneously, capturing aggregate USD risk.

\subsection{Strategy Construction}

\paragraph{DOLLAR.}
Go long an equal-weighted basket of all $N$ foreign currencies every month, unconditionally:
\begin{equation}
    R^{DOLLAR}_{t+1} = \frac{1}{N}\sum_{i=1}^N X^i_{t+1}.
\end{equation}
This is a permanent bet on USD depreciation. There is no signal: the investor simply holds
the foreign currency basket.

\paragraph{DOLLAR-CARRY.}
Go long (short) the same equal-weighted basket when the average interest rate differential
across currencies is positive (negative):
\begin{equation}
    R^{DOLLAR\text{-}CARRY}_{t+1} = \frac{\mathrm{sign}(\bar{\Delta r}_t)}{N}\sum_{i=1}^N X^i_{t+1},
    \qquad \bar{\Delta r}_t = \frac{1}{N}\sum_{i=1}^N \Delta r^i_t.
\end{equation}
The signal $\bar{\Delta r}_t$ is formed at month-end $t$ and applied to the following
month's returns, ensuring no look-ahead bias.

\subsection{Results --- Table 2 (Complete)}

Table~\ref{tab:table2} shows the complete six-strategy summary.

\paragraph{DOLLAR is essentially zero.}
DOLLAR earns a mean excess return of virtually zero ($-0.015\%$ annualised, $t = -0.01$):
over the full evaluation period, the USD neither appreciated nor depreciated significantly
against the equal-weighted G10 basket.

With an annualised volatility of $7.58\%$ and a Sharpe ratio of $-0.002$, the unconditional
long-basket strategy offered no compensation for its substantial directional risk.

\paragraph{DOLLAR-CARRY improves but remains insignificant.}
DOLLAR-CARRY earns $2.13\%$ annualised (Sharpe $= 0.28$), a meaningful improvement over the
unconditional strategy. The signal switches sides almost equally over the sample (179 months
long, 181 months short), confirming active conditioning.

However, its $t$-statistic of $1.54$ falls below the 1.96 threshold: the average cross-G10
rate differential is an imprecise predictor of broad USD direction at the monthly frequency.

\subsection{Extended Regression: Adding DOLLAR to Carry}

The extended regression in column~(2) of Table~\ref{tab:reg} adds DOLLAR as a third
regressor:
\begin{equation}
    R^{TS\text{-}CARRY}_{t+1} = \alpha + \beta_1 R^{CS\text{-}MOM}_{t+1}
        + \beta_2 R^{TS\text{-}MOM}_{t+1} + \beta_3 R^{DOLLAR}_{t+1} + u_{t+1}.
\end{equation}

\paragraph{$R^2$ jumps from 0.08\% to 16.1\%.}
The DOLLAR factor explains a substantial share of TS-CARRY variance: $\hat{\beta}_3 = 0.32$
with $t = 2.49$, the only significant slope coefficient in either regression.

This positive loading reflects a structural overlap. A sign-weighted carry strategy is, by
construction, persistently long on high-rate currencies (AUD, NZD, GBP) and short on the
lowest-rate currency (JPY). When the USD weakens globally, all foreign currencies tend to
appreciate simultaneously, mechanically lifting TS-CARRY returns.

\paragraph{The carry premium survives the dollar control.}
Despite the large $R^2$ increase, the carry alpha barely changes ($2.77\% \to 2.65\%$
annualised), and its $t$-statistic \emph{rises} from 2.57 to $\mathbf{3.02}$.

Controlling for dollar exposure removes variance from the carry return without reducing
its expected value, making the carry premium easier to detect statistically. The carry
premium is therefore genuine, structural, and distinct from dollar risk: it cannot be
explained as mere compensation for USD directional exposure.

%====================== SECTION 7: IN-SAMPLE PORTFOLIOS ======================
\section{In-sample Portfolio Analysis}
\label{sec:insample}

We now move from individual factors to portfolio construction. We have two sets of building
blocks:
\begin{itemize}
    \item the six basis currency excess returns $\{X^i_{t+1}\}_{i=1}^6$;
    \item the six factor strategies (CS-CARRY, TS-CARRY, CS-MOM, TS-MOM, DOLLAR,
          DOLLAR-CARRY).
\end{itemize}

For each set, we form three portfolios using different weighting schemes, then scale every
portfolio to a common annualised volatility of $10\%$ for direct Sharpe comparison.

\subsection{Weighting Schemes}

\paragraph{Equal-weight (EW).}
$w_i = 1/N$ for all $i$. Uses no signal at all and serves as a naive benchmark.

\paragraph{Risk-parity (RP).}
$w_i \propto 1/\hat{\sigma}_i$, where $\hat{\sigma}_i$ is the full in-sample standard
deviation of asset $i$. Each asset contributes equally to total portfolio risk; only the
diagonal of $\hat{\Sigma}$ is exploited.

\paragraph{Mean-variance (MV).}
$w^{*} \propto \hat{\Sigma}^{-1}\hat{\mu}$, where $\hat{\mu}$ and $\hat{\Sigma}$ are the
full in-sample sample mean and covariance matrix. This is the tangency portfolio that
maximises the in-sample Sharpe ratio.

\paragraph{Vol scaling.}
Each portfolio's raw return series is multiplied by a constant $k = 10\% / \hat{\sigma}_p$
so that the realised annualised volatility equals $10\%$. Since this is a constant
rescaling, it does not affect the Sharpe ratio --- it only allows mean returns to be
compared on a common risk basis.

\subsection{Results --- Table 3}

\begin{table}[H]
\centering
\caption{In-sample Portfolios --- Performance Statistics (Jan 1995 -- Dec 2024, scaled to 10\% vol)}
\label{tab:table3}
\renewcommand{\arraystretch}{1.3}
\begin{tabular}{lrrrrrr}
\toprule
 & \textbf{EW-CCY} & \textbf{RP-CCY} & \textbf{MV-CCY} & \textbf{EW-FAC} & \textbf{RP-FAC} & \textbf{MV-FAC} \\
\midrule
Mean (\%, ann.)  & $-$0.020 & $-$0.102 & 4.652 & 4.081 & 4.154 & \textbf{5.989} \\
$t$-statistic    & $-$0.011 & $-$0.056 & 2.548 & 2.235 & 2.275 & \textbf{3.280} \\
Std (\%, ann.)   & 10.000   & 10.000   & 10.000 & 10.000 & 10.000 & 10.000 \\
Sharpe ratio     & $-$0.002 & $-$0.010 & 0.465 & 0.408 & 0.415 & \textbf{0.599} \\
\bottomrule
\end{tabular}
\caption*{\footnotesize CCY = portfolios of the six basis currencies; FAC = portfolios of the
six factor strategies. All portfolios scaled ex-post to 10\% annualised volatility using
the full in-sample standard deviation.}
\end{table}

\begin{figure}[H]
    \centering
    \includegraphics[width=\textwidth]{fig_insample_portfolios.pdf}
    \caption{Cumulative returns of the six in-sample portfolios (10\%-vol scaled).
    MV-FAC dominates throughout; EW-CCY drifts horizontally around \$1, mirroring the
    near-zero DOLLAR return.}
    \label{fig:insample}
\end{figure}

\subsection{Discussion}

\subsubsection*{(c) RP-CCY vs MV-CCY: the value of correlations}

\textbf{MV-CCY (Sharpe 0.465) dominates RP-CCY (Sharpe $-0.010$) by a very wide margin.}

RP-CCY uses only the diagonal of the covariance matrix (individual volatilities), while
MV-CCY uses the full $\hat{\Sigma}$ including all pairwise correlations. The huge gap in
Sharpe ratios shows that the off-diagonal structure of the G10 covariance matrix carries
substantial information.

In particular, MV-CCY tilts heavily towards low-correlation combinations of currencies and
exploits the natural hedging relationships among G10 pairs --- information that RP-CCY
ignores entirely.

\subsubsection*{(d) MV-CCY vs MV-FAC: a theoretical paradox resolved by estimation noise}

The six factor strategies are themselves linear portfolios of the six basis currencies. The
factor space is therefore a \emph{linear subspace} of the currency portfolio space:
\[
    w_F^{\top} F = w_F^{\top} A X = (A^{\top} w_F)^{\top} X.
\]

Theoretically, with true population moments, the maximum Sharpe in the factor space cannot
exceed the maximum Sharpe in the currency space: $\mathrm{SR}^{*}_{\text{FAC}} \le
\mathrm{SR}^{*}_{\text{CCY}}$.

\textbf{In-sample, however, the opposite occurs: MV-FAC = 0.599 > MV-CCY = 0.465.}

The reversal is driven by estimation noise. The six factors were explicitly engineered to
have positive expected returns, which concentrates the signal in $\hat{\mu}_F$ and reduces
the noise in $\hat{\Sigma}_F^{-1}\hat{\mu}_F$ relative to the noisier $6\times 6$
currency-level system.

The factor space trades theoretical breadth for estimation precision --- and the trade-off
pays off in this sample.

\subsubsection*{(e) EW-CCY vs MV-FAC: a structural inequality}

\textbf{EW-CCY (Sharpe $-0.002$) is dominated by MV-FAC (Sharpe 0.599) by an order of
magnitude.}

This ordering is not coincidental: it is a mathematical necessity given the choice of
factor set. A key observation is that
\[
    \text{EW-CCY} \;\equiv\; \text{DOLLAR},
\]
since investing $1/N$ in each currency is exactly the equal-weighted basket return.
DOLLAR is one of the six factors used to build MV-FAC, so EW-CCY lies \emph{inside}
the factor span.

MV-FAC is by construction the tangency portfolio of the factor space, i.e.\ the maximum
in-sample Sharpe achievable with any combination of the six factors. Since EW-CCY is itself
a feasible factor combination (with all weight on DOLLAR), we always have
\[
    \mathrm{SR}^{\text{IS}}_{\text{MV-FAC}} \;\ge\; \mathrm{SR}^{\text{IS}}_{\text{EW-CCY}}.
\]

Therefore \textbf{it is not possible for EW-CCY to outperform MV-FAC in-sample as long as
DOLLAR (or any strategy that spans EW-CCY) is included in the factor set}.

The answer therefore depends explicitly on the factor set: if DOLLAR were removed, EW-CCY
would no longer lie in the factor span and could in principle outperform MV-FAC in some
specific sample.

%====================== SECTION 8: OUT-OF-SAMPLE ======================
\section{Out-of-sample Portfolio Analysis}
\label{sec:oos}

The in-sample portfolios of Section~\ref{sec:insample} use the full 360-month sample to
estimate moments. This look-ahead inflates their performance: a real-time investor can
only use data available at the formation date.

We now construct \emph{conditional mean-variance efficient} (CMVE) portfolios that satisfy
\begin{equation}
    w^{CMVE}_t = k_t \, \hat{\Sigma}_t^{-1} \hat{\mu}_t,
    \label{eq:cmve}
\end{equation}
where $\hat{\Sigma}_t$ and $\hat{\mu}_t$ are estimated using only data observed through
month $t$, and $k_t$ is a scaling factor that targets a constant ex-ante annualised
volatility of $10\%$.

\subsection{Conditional Covariance: EWMA}

The conditional covariance is estimated by an exponentially weighted moving average
(EWMA) with the standard RiskMetrics decay parameter $\lambda = 0.97$:
\begin{equation}
    \hat{\Sigma}_t = (1 - \lambda) \, X_t X_t^{\top} + \lambda \, \hat{\Sigma}_{t-1}.
\end{equation}

This recursion gives an effective half-life of approximately 23 months. It is initialised
at December 1994 with the sample covariance over the holdout window
(January 1991 -- December 1994, 48 observations).

\subsection{Conditional Mean: Two Approaches}

\paragraph{Expanding-window mean (Exp).}
A naive benchmark using all observations through $t$:
\begin{equation}
    \hat{\mu}^{\text{Exp}}_{i,t} = \frac{1}{N_t} \sum_{s=1}^{t} X^i_s.
\end{equation}

\paragraph{Fama-MacBeth mean (FM).}
Each month $s$, run a cross-sectional regression of the six excess returns on the carry
and momentum signals, with no intercept:
\begin{equation}
    X^i_{s+1} = \gamma_s \, \Delta r^i_s + \phi_s \, s^i_{s-12,s} + \varepsilon^i_s.
\end{equation}

Take the expanding-window average of the slope estimates,
\[
    \bar{\gamma}_t = \frac{1}{K_t} \sum_{s \le t-1} \hat{\gamma}_s,
    \qquad
    \bar{\phi}_t = \frac{1}{K_t} \sum_{s \le t-1} \hat{\phi}_s,
\]
and apply them to the current signals:
\begin{equation}
    \hat{\mu}^{\text{FM}}_{i,t} = \bar{\gamma}_t \, \Delta r^i_t + \bar{\phi}_t \, s^i_{t-12,t}.
\end{equation}

The FM regressions are pre-trained on the 36 cross-sections of the holdout window
(December 1991 -- November 1994), so the loop starts in January 1995 with non-degenerate
average coefficients.

\subsection{Vol Targeting}

For each formation date, the scaling factor is
\begin{equation}
    k_t = \frac{0.10}{\sqrt{\hat{\mu}_t^{\top} \hat{\Sigma}_t^{-1} \hat{\mu}_t \times 12}},
\end{equation}
which makes the ex-ante annualised volatility of $w^{CMVE}_t$ exactly equal to $10\%$
under the conditional moments.

Realised (ex-post) volatility may differ because the conditional moment estimates are
themselves noisy.

\subsection{Fama-MacBeth Coefficients --- Discussion (d)}

We accumulate 396 cross-sectional estimates of $(\hat{\gamma}_s, \hat{\phi}_s)$ from
December 1991 through November 2024.

\begin{table}[H]
\centering
\caption{Fama-MacBeth Coefficient Summary (1991--2024, 396 cross-sections)}
\label{tab:fmcoef}
\renewcommand{\arraystretch}{1.3}
\begin{tabular}{lrrrl}
\toprule
 & \textbf{Mean} & $t$-stat (OLS) & $t$-stat (NW, 12 lags) & Significance \\
\midrule
$\bar{\gamma}$ (carry slope)    & 1.739 & 2.50 & \textbf{2.89} & ✓ at 1\% \\
$\bar{\phi}$ (momentum slope)   & 0.016 & 1.18 & 1.43          & ✗ \\
\bottomrule
\end{tabular}
\caption*{\footnotesize OLS $t$-statistic uses the cross-section variability of monthly
slope estimates. Newey-West $t$-statistic (12 lags) is the standard HAC correction for
monthly financial data, allowing for serial correlation in the slopes.}
\end{table}

\paragraph{How the t-stats are computed.}
The cross-sectional regression is re-estimated each month, yielding a time series of slopes
$\{\hat{\gamma}_s, \hat{\phi}_s\}$ over $T = 396$ months.

The Fama-MacBeth $t$-statistic on $\bar{\gamma}$ is the time-series mean divided by its
standard error, optionally adjusted for autocorrelation via Newey-West with 12 lags.

\paragraph{Statistical significance.}
The carry slope is significant under both OLS and HAC inference (NW $t = 2.89$), confirming
a robust positive cross-sectional reward to high-rate currencies.

The momentum slope is not significant (NW $t = 1.43$), echoing the weak standalone
performance of CS-MOM and TS-MOM in Section~\ref{sec:momentum}.

\paragraph{Economic interpretation.}
A positive $\bar{\gamma}$ means that an extra unit of rate differential predicts higher
next-month excess return --- the carry premium and a violation of UIP.

A positive $\bar{\phi}$ (though small here) means that recently appreciating currencies
tend to keep appreciating --- the FX momentum premium.

\paragraph{Link to the factor strategies.}
The FM regression projects the six currency returns onto the \emph{exact same two signals}
used to build CS/TS-CARRY and CS/TS-MOM. The FM forecast $\hat{\mu}^{\text{FM}}_t$ is
therefore a linear combination of the carry and momentum signals, weighted by their
estimated cross-sectional risk premia.

The CCV-FM portfolio thus internalises the two strongest empirical regularities of the FX
cross-section in a single optimisation, rather than treating them as separate ad-hoc bets
--- which directly explains why CCV-FM dominates CCV-Exp out-of-sample.

\begin{figure}[H]
    \centering
    \includegraphics[width=\textwidth]{fig_fm_coefs.pdf}
    \caption{Evolution of the expanding-window FM coefficients $\bar{\gamma}_t$ (carry, top)
    and $\bar{\phi}_t$ (momentum, bottom) over the evaluation period. Both stabilise quickly
    after the first few years.}
    \label{fig:fmcoef}
\end{figure}

\subsection{CCV Portfolios --- Table 4}

\begin{table}[H]
\centering
\caption{Out-of-sample CMVE Portfolios vs.\ In-sample Benchmarks (Jan 1995 -- Dec 2024)}
\label{tab:table4}
\renewcommand{\arraystretch}{1.3}
\begin{tabular}{lrrrrr}
\toprule
 & \textbf{CCV-Exp} & \textbf{CCV-FM} & MV-CCY (IS) & MV-FAC (IS) & EW-CCY (IS) \\
 & (OOS) & (OOS) & & & \\
\midrule
Mean (\%, ann.)  & 2.127 & \textbf{6.182} & 4.652 & 5.989 & $-$0.020 \\
$t$-statistic    & 0.985 & \textbf{2.822} & 2.548 & 3.280 & $-$0.011 \\
Std (\%, ann.)   & 11.83 & 12.00          & 10.00 & 10.00 & 10.00 \\
Sharpe ratio     & 0.180 & \textbf{0.515} & 0.465 & 0.599 & $-$0.002 \\
\bottomrule
\end{tabular}
\caption*{\footnotesize CCV-Exp and CCV-FM are constructed in real time using only data
available at each formation date; their realised vol differs slightly from the 10\%
ex-ante target because of estimation noise. In-sample (IS) benchmarks reproduce columns
from Table~\ref{tab:table3} for direct comparison.}
\end{table}

\begin{figure}[H]
    \centering
    \includegraphics[width=\textwidth]{fig_oos_cumulative.pdf}
    \caption{Cumulative returns of the OOS portfolios (CCV-Exp, CCV-FM) alongside the
    in-sample MV benchmarks. CCV-FM tracks the in-sample MV-FAC closely despite operating
    in real time.}
    \label{fig:ooscum}
\end{figure}

\subsection{Performance Evaluation --- Discussion (f)}

\paragraph{CCV-FM dominates CCV-Exp.}
With a Sharpe of $0.515$ versus $0.180$, CCV-FM is roughly three times more efficient than
CCV-Exp. The FM forecast exploits the cross-sectional carry and momentum signals, which
concentrate the predictive content in a low-dimensional space.

The expanding sample mean used by CCV-Exp is overwhelmed by estimation noise, especially in
the early years where only a few dozen observations are available and the FX
signal-to-noise ratio is very low.

\paragraph{CCV-FM almost matches in-sample MV-FAC.}
The in-sample MV portfolios use the full 360-month sample to estimate both moments, which
is impossible in real time. CCV-FM, which forms portfolios using only past data, achieves
a Sharpe of $0.515$ --- within $0.08$ of the in-sample MV-FAC's $0.599$.

\textbf{This $0.08$ gap is the genuine, unavoidable cost of operating in real time with
this factor structure.}

\paragraph{Where the OOS degradation comes from.}
The dominant source of estimation error is the conditional mean $\hat{\mu}_t$. Second
moments converge much faster than means, so the EWMA covariance estimator is comparatively
well-behaved.

CCV-FM bypasses the worst part of the mean-estimation problem by imposing a parametric
structure (two factors: carry and momentum) instead of estimating six independent currency
means.

\paragraph{Realised vs ex-ante volatility.}
The scaling factor targets $10\%$ annualised vol using conditional moments. Realised
volatility is $11.83\%$ for CCV-Exp and $12.00\%$ for CCV-FM --- slightly above target due
to noise in real-time parameter estimates.

By construction, the in-sample portfolios hit exactly $10\%$ because they use ex-post
moments.

%====================== CONCLUSION ======================
\section{Conclusion}

This report documents the main characteristics of currency excess returns and a complete
factor-investing pipeline in the G10 FX market over 1995--2024.

\paragraph{Individual currencies: weak signals, high noise.}
No single currency delivers a Sharpe ratio above $0.20$ in absolute value or a significant
$t$-statistic. This motivates systematic multi-currency strategies rather than individual
bets.

\paragraph{Factor strategies: carry dominates.}
Among the six factor strategies, carry is the standout performer. CS-CARRY ($t = 2.17$)
and TS-CARRY ($t = 2.57$) both deliver statistically significant returns with Sharpe
ratios approaching $0.5$.

Momentum is much weaker (CS-MOM: $t = 0.15$, TS-MOM: $t = 1.29$, both insignificant) and
the unconditional DOLLAR factor earns essentially zero. DOLLAR-CARRY improves on DOLLAR
but does not reach significance.

\paragraph{The carry premium is robust.}
The carry premium survives controls for both momentum and dollar exposure: its alpha
$t$-statistic actually \emph{rises} from 2.57 standalone to 3.02 after controlling for
MOM + DOLLAR. The premium is therefore not an artefact of momentum or USD risk.

\paragraph{In-sample portfolios: factors beat raw currencies.}
Of the six 10\%-vol-scaled portfolios, MV-FAC achieves the highest Sharpe ($0.599$),
followed by MV-CCY ($0.465$). Naive equal-weight currency portfolios are essentially
unprofitable.

The factor space achieves better in-sample performance precisely because the factors are
designed to concentrate the cross-sectional signal.

\paragraph{Out-of-sample CMVE: parsimony wins.}
The Fama-MacBeth-based CCV-FM achieves a Sharpe of $0.515$ in real time --- roughly three
times the naive expanding-mean CCV-Exp ($0.180$), and within $0.08$ of the in-sample
MV-FAC.

This confirms that the practical value of mean-variance optimisation in FX hinges critically
on imposing a parsimonious structure (two factors: carry and momentum) on the conditional
mean. Estimating six independent means from monthly FX data is hopeless; estimating two
risk premia from a cross-sectional regression is feasible.

\paragraph{Synthesis.}
Carry, momentum, and dollar capture distinct sources of FX risk premia, with carry being
the strongest and momentum the weakest. Their combination through a Fama-MacBeth
cross-sectional projection delivers a real-time Sharpe of $0.5$, demonstrating that the
asset-pricing intuition translates directly into implementable portfolio gains.

\end{document}
